
\chapter{Boundary Value Problems for Second Order Ordinary
Differential Equations}


\noindent
IN THIS CHAPTER we discuss  boundary value problems and eigenvalue
problems for linear second order ordinary differential equations.

\medskip\noindent
Section~13.1 discusses point two-point boundary value
problems for linear second order ordinary differential equations.


\medskip\noindent
Section~13.2 deals with generalizations of the eigenvalue
problems considered in Section~11.1

\vspace*{4in}
\centerline{677}

\newpage

\newsection{1}
{Boundary Value Problems for Second Order Ordinary
Differential Equations}
{Two-Point Boundary Value Problems}
\subpdfbookmark{Section 13.1 Two-Point Boundary Value Problems} {section:13.1}
\vskip14pt


\renewcommand{\thissection}
{\sectiontitle{TWO-POINT BOUNDARY VALUE PROBLEMS}}
\thissection



\noindent
In Section~5.3 we considered initial value problems for the
linear second order equation
\begin{equation} \label{eq:13.1.1}
P_{0}(x)y''+P_{1}(x)y'+P_{2}(x)y=F(x).
\end{equation}
Suppose $P_{0}$, $P_{1}$, $P_{2}$, and $F$ are continuous and $P_{0}$
has no zeros
on an
open interval $(a,b)$. From Theorem~\ref{thmtype:5.3.1}, if $x_{0}$ is in
$(a,b)$ and  $k_{1}$ and $k_{2}$ are arbitrary real numbers then
\eqref{eq:13.1.1} has a unique solution on $(a,b)$ such that $y(x_{0})=k_{1}$
and $y'(x_{0})=k_{2}$. Now we consider a different problem for
\eqref{eq:13.1.1}.

{\bf PROBLEM}\quad
Suppose $P_{0}$, $P_{1}$, $P_{2}$, and $F$ are  continous and  $P_{0}$
has no zeros on a closed interval $[a,b]$. Let $\alpha$, $\beta$,
$\rho$, and $\delta$ be real numbers such that
\begin{equation} \label{eq:13.1.2}
\alpha^{2}+\beta^{2}\ne 0\text{\;    and\; \;  } \rho^{2}+\delta^{2}\ne 0,
\end{equation}
and let $k_{1}$ and $k_{2}$ be arbitrary real numbers. Find a solution of
\begin{equation} \label{eq:13.1.3}
P_{0}(x)y''+P_{1}(x)y'+P_{2}(x)y=F(x)
\end{equation}
on the closed interval $[a,b]$ such that
\begin{equation} \label{eq:13.1.4}
\alpha y(a)+\beta y'(a)=k_{1}
\end{equation}
and
\begin{equation} \label{eq:13.1.5}
\rho y(b)+\delta y'(b)=k_{2}.
\end{equation}


\medskip
The assumptions  stated in this problem apply throughout this section and
won't be repeated. Note that we imposed conditions on $P_{0}$,
$P_{1}$, $P_{2}$, and $F$  on the \emph{closed} interval $[a,b]$, and we
are interested in solutions of \eqref{eq:13.1.3} on the closed interval.
This is different from the situation considered in Chapter~5,
where we imposed conditions on $P_{0}$, $P_{1}$, $P_{2}$, and $F$  on the
\emph{open} interval $(a,b)$ and we were interested in solutions on the
open interval. There
is really no problem here;   we can always extend  $P_{0}$, $P_{1}$,
$P_{2}$,
and $F$ to an open interval $(c,d)$ (for example, by defining them to be
constant on $(c,d]$ and $[b,d)$) so that they are continuous and $P_{0}$
has no zeros on $[c,d]$. Then we can apply the theorems from
Chapter~5 to the equation
$$
y''+\frac{P_{1}(x)}{P_{0}(x)}y'+
\frac{P_{2}(x)}{P_{0}(x)}y=\frac{F(x)}{P_{0}(x)}
$$
on $(c,d)$ to draw conclusions about solutions of \eqref{eq:13.1.3} on $[a,b]$.

We call $a$ and $b$ \emph{boundary points.}  The conditions \eqref{eq:13.1.4}
and \eqref{eq:13.1.5}  are \emph{boundary conditions}, and the problem is a
\emph{two-point boundary value problem} or, for simplicity, a
\emph{boundary value problem}. (We used  similar terminology in
Chapter~12 with a different meaning; both meanings are in common usage.)
  We require \eqref{eq:13.1.2} to insure that we're imposing a sensible
condition at each boundary point. For example, if $\alpha^{2}+\beta^{2}=0$
then $\alpha=\beta=0$, so $\alpha y(a)+\beta y'(a)=0$ for all choices of
$y(a)$ and $y'(a)$. Therefore  \eqref{eq:13.1.4} is an impossible condition if
$k_{1}\ne0$, or no condition at all if $k_{1}=0$.

We abbreviate \eqref{eq:13.1.1} as $Ly=F$, where
$$
Ly=P_{0}(x)y''+P_{1}(x)y'+P_{0}(x)y,
$$
and we denote
$$
B_{1}(y)=\alpha y(a)+\beta y'(a)
\text{\;    and\; \;  }
B_{2}(y)=\rho y(b)+\delta y'(b).
$$
We combine \eqref{eq:13.1.3}, \eqref{eq:13.1.4}, and \eqref{eq:13.1.5} as
\begin{equation} \label{eq:13.1.6}
Ly=F,\quad \quad B_{1}(y)=k_{1},\quad \quad B_{2}(y)=k_{2}.
\end{equation}
This boundary value problem is \emph{homogeneous} if $F=0$
and $k_{1}=k_{2}=0$;   otherwise it's \emph{nonhomogeneous}.

We leave it to you (Exercise~\ref{exer:13.1.1}) to verify that $B_{1}$
and $B_{2}$ are linear operators; that is, if $c_{1}$  and $c_{2}$
are constants then
\begin{equation} \label{eq:13.1.7}
B_{i}(c_{1}y_{1}+c_{2}y_{2})=c_{1}B_{i}(y_{1})+c_{2}B_{i}(y_{2}),\quad
i=1,2.
\end{equation}
The next three examples show that the question of existence and uniqueness
for solutions of boundary value problems is more complicated than
 for initial value problems.

\begin{example}\label{example:13.1.1}    \rm
Consider the boundary value problem
$$
y''+y=1,\quad y(0)=0, \quad y(\pi/2)=0.
$$
The general solution of $y''+y=1$ is
$$
y=1+c_{1}\sin x+c_{2} \cos x,
$$
so $y(0)=0$ if  and only if $c_{2}=-1$ and $y(\pi/2)=0$
if and only if $c_{1}=-1$. Therefore
$$
 y =1-\sin x -\cos x
$$
is the unique solution of the boundary value problem.
\end{example}

\begin{example}\label{example:13.1.2}    \rm
Consider the boundary value problem
$$
y''+y=1,\quad y(0)=0, \quad y(\pi)=0.
$$
Again, the general solution of $y''+y=1$ is
$$
y=1+c_{1}\sin x+c_{2} \cos x,
$$
so $y(0)=0$ if  and only if $c_{2}=-1$, but $y(\pi)=0$ if and only if
$c_{2}=1$. Therefore the boundary value problem has no solution.
\end{example}

\begin{example}\label{example:13.1.3}    \rm
Consider the boundary value problem
$$
y''+y=\sin 2x,\quad y(0)=0, \quad y(\pi)=0.
$$
You can use the method of undetermined coefficients
(Section~5.5) to find that the general solution of
$y''+y=\sin 2x$  is
$$
y=-\frac{\sin2x}{3}+c_{1}\sin x+c_{2}\cos x.
$$
The boundary conditions $y(0)=0$ and $y(\pi)=0$ both require that
$c_{2}=0$, but they don't restrict $c_{1}$. Therefore the boundary value
problem has infinitely many solutions
$$
y=-\frac{\sin2x}{3}+c_{1}\sin x,
$$
where $c_{1}$  is arbitrary.
\end{example}


\begin{theorem}\color{blue} \label{thmtype:13.1.1}
If $z_{1}$ and $z_{2}$ are solutions of $Ly=0$  such that either
$B_{1}(z_{1})=B_{1}(z_{2})=0$ or
$B_{2}(z_{1})=B_{2}(z_{2})=0,$  then $\{z_{1},z_{2}\}$  is linearly
dependent. Equivalently$,$ if\,  $\{z_{1},z_{2}\}$ is linearly
independent, then
$$
B_{1}^{2}(z_{1})+B_{1}^{2}(z_{2})\ne 0 \text{\;    and\; \;  }
B_{2}^{2}(z_{1})+B_{2}^{2}(z_{2})\ne 0.
$$
\end{theorem}

\medskip
\proof
Recall that $B_{1}(z)=\alpha z(a)+\beta z'(a)$ and
$\alpha^{2}+\beta^{2}\ne0$.
Therefore, if $B_{1}(z_{1})=B_{1}(z_{2})=0$ then $(\alpha,\beta)$
is a nontrivial solution of  the system
\begin{eqnarray*}
\alpha z_{1}(a)+\beta z_{1}'(a)&=& 0\\
\alpha z_{2}(a)+\beta z_{}'(a)&=& 0.
\end{eqnarray*}
This  implies that
$$
z_{1}(a)z_{2}'(a)-z_{1}'(a)z_{2}(a)=0,
$$
so $\{z_{1},z_{2}\}$ is linearly dependent, by
Theorem~\ref{thmtype:5.1.6}. We leave it to you to show
that $\{z_{1},z_{2}\}$ is linearly dependent   if
$B_{2}(z_{1})=B_{2}(z_{2})=0$.


\begin{theorem}\color{blue} \label{thmtype:13.1.2}
The following statements are equivalent$;$ that is$,$ they are either
all true or all false$.$
\begin{alist}
\item    %(a)
 There's a fundamental set $\{z_{1},z_{2}\}$ of solutions of
$Ly=0$ such that
\begin{equation} \label{eq:13.1.8}
B_{1}(z_{1})B_{2}(z_{2})-B_{1}(z_{2})B_{2}(z_{1})\ne0.
\end{equation}
\item %(b)
If  $\{y_{1},y_{2}\}$ is a fundamental set of solutions of $Ly=0$  then
\begin{equation} \label{eq:13.1.9}
B_{1}(y_{1})B_{2}(y_{2})-B_{1}(y_{2})B_{2}(y_{1})\ne0.
\end{equation}
\item % (c)
For each continuous $F$ and pair of constants $(k_{1},k_{2}),$  the
boundary value problem
$$
Ly=F, \quad B_{1}(y)=k_{1},\quad B_{2}(y)=k_{2}
$$
has a unique solution$.$
\item % (d)
The homogeneous boundary value problem
\begin{equation} \label{eq:13.1.10}
Ly=0,\quad  B_{1}(y)=0, \quad B_{2}(y)=0
\end{equation}
has only the trivial solution $y=0$.
\item %(e)
The homogeneous equation $Ly=0$ has linearly independent solutions
$z_{1}$ and $z_{2}$ such that $B_{1}(z_{1})=0$ and $B_{2}(z_{2})=0.$
\end{alist}
\end{theorem}



\proof
We'll show  that
$$
\part{a}\implies
\part{b}\implies
\part{c}\implies
\part{d}\implies
\part{e}\implies
\part{a}.
$$

\part{a} $\implies$ \part{b}:
Since $\{z_{1},z_{2}\}$ is a fundamental set of solutions for $Ly=0$,
there are
constants
$a_{1}$, $a_{2}$, $b_{1}$, and $b_{2}$  such that
\begin{equation} \label{eq:13.1.11}
\begin{aligned}
y_{1} &=a_{1}z_{1}+a_{2}z_{2}\\
y_{2} &=b_{1}z_{1}+b_{2}z_{2}.
\end{aligned}
\end{equation}
Moreover,
\begin{equation} \label{eq:13.1.12}
\left|\begin{array}{ccccccc}
a_{1}&a_{2}\\b_{1}&b_{2}
\end{array}\right|\ne0.
\end{equation}
because if this determinant were zero, its rows would be linearly
dependent and therefore
$\{y_{1},y_{2}\}$ would be linearly dependent, contrary to our assumption
that  $\{y_{1},y_{2}\}$  is a fundamental set of solutions of $Ly=0$.
From \eqref{eq:13.1.7} and \eqref{eq:13.1.11},
$$
\left[\begin{array}{ccccccc}
B_{1}(y_{1})&B_{2}(y_{1})\\B_{1}(y_{2})&B_{2}(y_{2})
\end{array}\right]=
\left[\begin{array}{ccccccc}
a_{1}&a_{2}\\b_{1}&b_{2}
\end{array}\right]
\left[\begin{array}{ccccccc}
B_{1}(z_{1})&B_{2}(z_{1})\\
B_{1}(z_{2})&B_{2}(z_{2})
\end{array}\right].
$$
Since the determinant of a product of matrices is the product of the
determinants of the matrices, \eqref{eq:13.1.8} and \eqref{eq:13.1.12} imply
\eqref{eq:13.1.9}.


\part{b} $\implies$ \part{c}:
Since $\{y_{1},y_{2}\}$ is a fundamental set of solutions of
$Ly=0$, the general solution of  $Ly=F$ is
$$
y=y_{p}+c_{1}y_{1}+c_{2}y_{2},
$$
where $c_{1}$ and $c_{2}$ are arbitrary constants and  $y_{p}$ is a
particular solution of
$Ly=F$. To satisfy the boundary conditions, we must choose $c_{1}$ and
$c_{2}$ so that
\begin{eqnarray*}
k_{1}&=&B_{1}(y_{p})+c_{1}B_{1}(y_{1})+c_{2}B_{1}(y_{2})\\
k_{2}&=&B_{2}(y_{p})+c_{1}B_{2}(y_{1})+c_{2}B_{2}(y_{2}),
\end{eqnarray*}
(recall \eqref{eq:13.1.7}), which is equivalent to
$$
\begin{aligned}
c_{1}B_{1}(y_{1})+c_{2}B_{1}(y_{2}) &= k_{1}-B_{1}(y_{p}) \\
c_{1}B_{2}(y_{1})+c_{2}B_{2}(2_{2}) &= k_{2}-B_{2}(y_{p}).
\end{aligned}
$$
From  \eqref{eq:13.1.9}, this system always has a  unique solution
 $(c_{1},c_{2})$.

\part {c} $\implies$ \part{d}: Obviously, $y=0$ is a solution of
\eqref{eq:13.1.10}. From \part{c} with $F=0$ and $k_{1}=k_{2}=0$, it's the only
solution.

\part{d} $\implies$ \part{e}: Let  $\{y_{1},y_{2}\}$ be a fundamental
system for $Ly=0$ and  let
$$
z_{1}=B_{1}(y_{2})y_{1}-B_{1}(y_{1})y_{2}\text{\;    and\; \;  }
z_{2}=B_{2}(y_{2})y_{1}-B_{2}(y_{1})y_{2}.
$$
Then $B_{1}(z_{1})=0$  and $B_{2}(z_{2})=0$. To see
that $z_{1}$ and $z_{2}$ are linearly independent, note that
\begin{eqnarray*}
a_{1}z_{1}+a_{2}z_{2}&=&
a_{1}[B_{1}(y_{2})y_{1}-B_{1}(y_{1})y_{2}]+
a_{2}[B_{2}(y_{2})y_{1}-B_{2}(y_{1})y_{2}]\\
&=&
[B_{1}(y_{2})a_{1}+B_{2}(y_{2})a_{2}]y_{1}-[B_{1}(y_{1})a_{1}+B_{2}(y_{1})a_{2}]y_{2}.
\end{eqnarray*}
Therefore, since $y_{1}$ and $y_{2}$ are linearly independent,
$a_{1}z_{1}+a_{2}z_{2}=0$ if and only if
$$
\left[\begin{array}{ccccccc}
B_{1}(y_{1})&B_{2}(y_{1})\\
B_{1}(y_{2})&B_{2}(y_{2})
\end{array}\right]
\left[\begin{array}{ccccccc}
a_{1}\\a_{2}
\end{array}\right]=
\left[\begin{array}{ccccccc}
0\\0
\end{array}\right].
$$
If this system has a nontrivial solution then so does the system
$$
\left[\begin{array}{ccccccc}
B_{1}(y_{1})&B_{1}(y_{2})\\
B_{2}(y_{1})&B_{2}(y_{2})
\end{array}\right]
\left[\begin{array}{ccccccc}
c_{1} \\c_{2}
\end{array}\right]=
\left[\begin{array}{ccccccc}
0\\0
\end{array}\right].
$$
This and \eqref{eq:13.1.7} imply that $y=c_{1}z_{1}+c_{2}z_{2}$ is
a nontrivial solution of \eqref{eq:13.1.10}, which contradicts~\part{d}.


\part{e} $\implies$ \part{a}.
Theorem~\ref{thmtype:13.1.1} implies that if  $B_{1}(z_{1})=0$  and
$B_{2}(z_{2})=0$ then $B_{1}(z_{2})\ne0$ and $B_{2}(z_{1})\ne0$.
This implies~\eqref{eq:13.1.8}, which completes the proof.



\begin{example}\label{example:13.1.4} \rm
Solve the boundary value problem
\begin{equation} \label{eq:13.1.13}
x^{2}y''-2xy'+2y-2x^{3}=0,\quad y(1)=4,\quad y'(2)=3,
\end{equation}
given that $\{x,x^{2}\}$ is a fundamental set of solutions
of the complementary equation
\end{example}

\solution
Using variation of parameters (Section~5.7), you can show
that
$y_{p}=x^{3}$ is a solution of the complementary equation
$$
x^{2}y''-2xy'+2y=2x^{3}=0.
$$
Therefore the solution of
\eqref{eq:13.1.13} can be written as
$$
y=x^{3}+c_{1}x+c_{2}x^{2}.
$$
Then
$$
y'=3x^{2}+c_{1}+2c_{2}x,
$$
and imposing the boundary conditions yields the system
\begin{eqnarray*}
c_{1}+\phantom{4}c_{2}&=&\phantom{-}3\\ c_{1}+4c_{2}&=&-9,
\end{eqnarray*}
so $c_{1}=7$ and $c_{2}=-4$.
Therefore
$$
y=x^{3}+7x-4x^{2}
$$
is the unique solution of \eqref{eq:13.1.13}

\begin{example}\label{example:13.1.5} \rm
Solve the boundary value problem
$$
y''-7y'+12y=4e^{2x},\quad y(0)=3,\quad y(1)=5e^{2}.
$$
\end{example}

\solution
From Example~\ref{example:5.4.1}, $y_{p}=2e^{2x}$ is a particular solution
of
\begin{equation} \label{eq:13.1.14}
y''-7y'+12y=4e^{2x}.
\end{equation}
Since $\{e^{3x},e^{4x}\}$ is a fundamental set for the complementary
equation, we could write the solution of \eqref{eq:13.1.13} as
$$
y=2e^{2x}+c_{1}e^{3x}+c_{2}e^{4x}
$$
and determine $c_{1}$ and $c_{2}$ by imposing the boundary conditions.
However, this would lead to some tedious algebra, and the form of the
solution would be very unappealing. (Try it!) In this case it's convenient
to use the fundamental system $\{z_{1},z_{2}\}$ mentioned in
Theorem~\ref{thmtype:13.1.2}\part{e}; that is, we choose
$\{z_{1},z_{2}\}$ so that $B_{1}(z_{1})=z_{1}(0)=0$ and
$B_{2}(z_{2})=z_{2}(1)=0$. It is easy to see that
$$
z_{1}=e^{3x}-e^{4x}\text{\;    and\; \;  }z_{2}=e^{3(x-1)}-e^{4(x-1)}
$$
satisfy these requirements.
Now we write the solution of \eqref{eq:13.1.14} as
$$
y=2e^{2x}+c_{1}\left(e^{3x}-e^{4x}\right)+c_{2}\left(e^{3(x-1)}-e^{4(x-1)}\right).
$$
Imposing the boundary conditions $y(0)=3$  and $y(1)=5e^{2}$ yields
$$
3=2+c_{2}e^{-4}(e-1)\text{\;    and\; \;  }
5e^{2}=2e^{2}+c_{1}e^{3}(1-e).
$$
Therefore
$$
c_{1}=\frac{3}{e(1-e)},\quad c_{2}=\frac{e^{4}}{e-1},
$$
and
$$
y=2e^{2x}+\frac{3}{e(1-e)}(e^{3x}-e^{4x})+
\frac{e^{4}}{e-1}(e^{3(x-1)}-e^{4(x-1)}).
$$


Sometimes  it's useful to have a formula for the solution of a general
boundary problem. Our next theorem addresses this question.

\begin{theorem}\color{blue} \label{thmtype:13.1.3}
Suppose the homogeneous boundary value problem
\begin{equation} \label{eq:13.1.15}
Ly=0,\quad B_{1}(y)=0,\quad B_{2}(y)=0
\end{equation}
has only the trivial solution$.$ Let $y_{1}$ and $y_{2}$ be linearly
independent solutions   of $Ly=0$ such that $B_{1}(y_{1})=0$ and
$B_{2}(y_{2})=0,$ and let
$$
W=y_{1}y_{2}'-y_{1}'y_{2}.
$$
Then the unique solution of
\begin{equation} \label{eq:13.1.16}
Ly=F,\quad B_{1}(y)=0,\quad B_{2}(y)=0
\end{equation}
is
\begin{equation} \label{eq:13.1.17}
y(x)= y_{1}(x)\int_{x}^{b}\frac{F(t)y_{2}(t)}{P_{0}(t)W(t)}\,dt+
y_{2}(x)\int_{a}^{x}\frac{F(t)y_{1}(t)}{P_{0}(t)W(t)}\,dt.
\end{equation}
\end{theorem}

\medskip
\proof In Section~5.7  we saw that if
\begin{equation} \label{eq:13.1.18}
y=u_{1}y_{1}+u_{2}y_{2}
\end{equation}
where
\begin{eqnarray*}
u_{1}'y_{1}+u_{2}'y_{2}&=&0\\
u_{1}'y_{1}'+u_{2}'y_{2}'&=&F,
\end{eqnarray*}
then  $Ly=F$. Solving for $u_{1}'$ and $u_{2}'$ yields
$$
u_{1}'=-\frac{Fy_{2}}{P_{0}W}\text{\;    and\; \;  }
u_{2}'=\frac{Fy_{1}}{P_{0}W},
$$
which hold if
$$
u_{1}(x)=\int_{x}^{b}\frac{F(t)y_{2}(t)}{P_{0}(t)W(t)}\,dt
\text{\;    and\; \;  }
u_{2}(x)=\int_{a}^{x}\frac{F(t)y_{1}(t)}{P_{0}(t)W(t)}\,dt.
$$
This and \eqref{eq:13.1.18} show that \eqref{eq:13.1.17} is a solution of $Ly=F$.
Differentiating \eqref{eq:13.1.17}  yields
\begin{equation} \label{eq:13.1.19}
y'(x)=y_{1}'(x)\int_{x}^{b}\frac{F(t)y_{2}(t)}{P_{0}(t)W(t)}\,dt+
y_{2}'(x)\int_{a}^{x}\frac{F(t)y_{1}(t)}{P_{0}(t)W(t)}\,dt.
\end{equation}
(Verify.) From \eqref{eq:13.1.17} and \eqref{eq:13.1.19},
$$
B_{1}(y)=B_{1}(y_{1})\int_{a}^{b}\frac{F(t)y_{2}(t)}{P_{0}(t)W(t)}\,dt=0
$$
because $B_{1}(y_{1})=0$, and
$$
B_{2}(y)=B_{2}(y_{2})\int_{a}^{b}\frac{F(t)y_{1}(t)}{P_{0}(t)W(t)}\,dt=0
$$
because $B_{2}(y_{2})=0$. Hence, $y$  satisfies \eqref{eq:13.1.16}.
This completes the proof.

We  can rewrite \eqref{eq:13.1.17} as
\begin{equation} \label{eq:13.1.20}
y=\int_{a}^{b} G(x,t)F(t)\,dt,
\end{equation}
where
$$
G(x,t)=
\begin{cases}
\dst{\frac{y_{1}(t)y_{2}(x)}{P_{0}(t)W(t)}}&,\quad a\le t\le x,\\[.1in]
\dst{\frac{y_{1}(x)y_{2}(t)}{P_{0}(t)W(t)}} &,\quad x\le t\le b.
\end{cases}.
$$

This is the   \emph{Green's function} for the boundary value problem
\eqref{eq:13.1.16}.
The Green's function  is related to the boundary value problem
\eqref{eq:13.1.16} in much the
same way that the inverse of a square matrix ${\mathbf A}$ is related to
the linear algebraic system ${\mathbf y}={\mathbf A}{\mathbf x}$;   just as
we substitute the given vector ${\mathbf y}$ into the formula
${\mathbf x}={\mathbf A}^{-1}{\mathbf y}$ to solve
 ${\mathbf y} ={\mathbf Ax}$, we substitute the given
function $F$ into the formula \eqref{eq:13.1.20} to obtain the solution of
\eqref{eq:13.1.16}. The analogy goes further:
just as ${\mathbf A}^{-1}$ exists if and only if
 ${\mathbf A}{\mathbf x} ={\mathbf 0}$
has only the trivial solution, the boundary value problem  \eqref{eq:13.1.16}
has a Green's function if and only the homogeneous boundary value problem
\eqref{eq:13.1.15}  has only the trivial solution.

We leave it to you (Exercise~\ref{exer:13.1.32}) to show that the
assumptions of Theorem~\ref{thmtype:13.1.3} imply that the unique solution
of the boundary value problem
$$
Ly=F,\quad B_{1}(y)=k_{1},\quad B_{2}(y)=k_{2}
$$
is
$$
y(x)=\int_{a}^{b}G(x,t)F(t)\,dt +\frac{k_{2}}{B_{2}(y_{1})}y_{1}+
\frac{k_{1}}{B_{1}(y_{2})}y_{2}.
$$

\begin{example}\label{example:13.1.6}\rm
Solve the boundary value problem
\begin{equation} \label{eq:13.1.21}
y''+y=F(x).\quad y(0)+y'(0)=0,\quad y(\pi)-y'(\pi)=0,
\end{equation}
and find the Green's function for this problem.
\end{example}

\solution
Here
$$
B_{1}(y)=y(0)+y'(0)   \text{\;    and\; \;  } B_{2}(y)=y(\pi)-y'(\pi).
$$
Let $\{z_{1},z_{2}\}=\{\cos x,\sin x\}$, which is a fundamental set of
solutions of $y''+y=0$. Then
\begin{eqnarray*}
B_{1}(z_{1})&=&(\cos x-\sin x)\big|_{x=0}=\phantom{-}1\\
B_{2}(z_{1})&=&(\cos x+\sin x)\big|_{x=\pi}=-1
\end{eqnarray*}
and
\begin{eqnarray*}
B_{1}(z_{2})&=&(\sin x+\cos x)\big|_{x=0}=1\\
B_{2}(z_{2})&=&(\sin x-\cos x)\big|_{x=\pi}=1.
\end{eqnarray*}
Therefore
$$
B_{1}(z_{1})B_{2}(z_{2})-B_{1}(z_{2})B_{2}(z_{1})=2,
$$
so Theorem~\ref{thmtype:13.1.2} implies that \eqref{eq:13.1.21} has a unique
solution. Let
$$
y_{1}=B_{1}(z_{2})z_{1}-B_{1}(z_{1})z_{2}= \cos x-\sin x
$$
 and
$$
y_{2}=B_{2}(z_{2})z_{1}-B_{2}(z_{1})z_{2}=  \cos x+\sin x.
$$
Then $B_{1}(y_{1})=0$, $B_{2}(y_{2})=0$, and the Wronskian of
$\{y_{1},y_{2}\}$ is
$$
W(x)=
\left|\begin{array}{rrccccc}
\cos x-\sin x& \cos x+ \sin x\\
-\sin x-\cos x&-\sin x+\cos x
\end{array}\right|=2.
$$
Since $P_{0}=1$, \eqref{eq:13.1.17} yields the solution
\begin{eqnarray*}
y(x)&=&\dst{\frac{\cos x-\sin x}{2}}\int_{x}^{\pi}F(t)(\cos t+\sin t)\,dt\\
&&+\dst{\frac{\cos x +\sin x}{2}}\int_{0}^{x}F(t)(\cos t-\sin t)\,dt.
\end{eqnarray*}
The Green's function is
$$
G(x,t)=
\begin{cases}
\dst{\frac{(\cos t-\sin t)(\cos x+\sin x)}{2}},&\quad  0\le t\le
x,\\[.1in]
\dst{\frac{(\cos x-\sin x)(\cos t+\sin t)}{2}},&\quad x\le t\le \pi.
\end{cases}
$$

We'll now consider the situation not covered by Theorem~\ref{thmtype:13.1.3}.

\begin{theorem}\color{blue} \label{thmtype:13.1.4}
Suppose the homogeneous boundary value problem
\begin{equation} \label{eq:13.1.22}
Ly=0,\quad B_{1}(y)=0,\quad B_{2}(y)=0
\end{equation}
has a nontrivial solution $y_{1},$  and let $y_{2}$ be  any solution of
$Ly=0$ that isn't  a constant multiple of $y_{1}.$  Let
$W=y_{1}y_{2}'-y_{1}'y_{2}.$ If
\begin{equation} \label{eq:13.1.23}
\int_{a}^{b}\frac{F(t)y_{1}(t)}{P_{0}(t)W(t)}\,dt=0,
\end{equation}
then the homogeneous boundary value problem
\begin{equation} \label{eq:13.1.24}
Ly=F, \quad B_{1}(y)=0,\quad B_{2}(y)=0
\end{equation}
has infinitely many solutions$,$ all of the form $y=y_{p}+c_{1}y_{1},$
where
$$
y_{p}=y_{1}(x)
\int_{x}^{b}\frac{F(t)y_{2}(t)}{P_{0}(t)W(t)}\,dt+
y_{2}(x)
\int_{a}^{x}\frac{F(t)y_{1}(t)}{P_{0}(t)W(t)}\,dt
$$
and $c_{1}$ is a constant$.$ If
$$
\int_{a}^{b}\frac{F(t)y_{1}(t)}{P_{0}(t)W(t)}\,dt \ne0,
$$
then \eqref{eq:13.1.24}  has no solution$.$
\end{theorem}

\proof
From the proof of Theorem~\ref{thmtype:13.1.3}, $y_{p}$ is a particular solution
of $Ly=F$,
and
$$
y_{p}'(x)=y_{1}'(x)\int_{x}^{b}\frac{F(t)y_{2}(t)}{P_{0}(t)W(t)}\,dt
+y_{2}'(x)\int_{a}^{x}\frac{F(t)y_{1}(t)}{P_{0}(t)W(t)}\,dt.
$$
Therefore the general solution of \eqref{eq:13.1.22} is of the form
$$
y=y_{p}+c_{1}y_{1}+c_{2}y_{2},
$$
where $c_{1}$ and $c_{2}$  are constants.  Then
\begin{eqnarray*}
B_{1}(y)&=&
B_{1}(y_{p}+c_{1}y_{1}+c_{2}y_{2})=
B_{1}(y_{p})+c_{1}B_{1}(y_{1})+c_{2}B_{1}y_{2}\\
&=&B_{1}(y_{1})\int_{a}^{b}\frac{F(t)y_{2}(t)}{P_{0}(t)W(t)}\,dt
+c_{1}B_{1}(y_{1})+c_{2}B_{1}(y_{2})\\
&=&c_{2}B_{1}(y_{2})
\end{eqnarray*}
Since $B_{1}(y_{1})=0$, Theorem~\ref{thmtype:13.1.1} implies that
$B_{1}(y_{2})\ne0$; hence, $B_{1}(y)=0$ if and only if $c_{2}=0$.
Therefore $y=y_{p}+c_{1}y_{1}$ and
\begin{eqnarray*}
B_{2}(y)&=&B_{2}(y_{p}+c_{1}y_{1})
=B_{2}(y_{2})\int_{a}^{b}\frac{F(t)y_{1}(t)}{P_{0}(t)W(t)}\,dt+c_{1}B_{2}(y_{1})\\
&=&B_{2}(y_{2})\int_{a}^{b}\frac{F(t)y_{1}(t)}{P_{0}(t)W(t)}\,dt,
\end{eqnarray*}
since $B_{2}(y_{1})=0$.  From Theorem~\ref{thmtype:13.1.1},
$B_{2}(y_{2})\ne0$ (since $B_{2}(y_{1}=0$).
Therefore $Ly=0$   if and only if  \eqref{eq:13.1.23} holds.
This completes the proof.




\begin{example}\label{example:13.1.7} \rm
Applying Theorem~\ref{thmtype:13.1.4} to the boundary value problem
\begin{equation} \label{eq:13.1.25}
y''+y=F(x),\quad y(0)=0,\quad y(\pi)=0
\end{equation}
explains the  Examples~\ref{example:13.1.2} and \ref{example:13.1.3}. The
complementary equation $y''+y=0$ has the linear independent solutions
$y_{1}=\sin x$ and $y_{2}=\cos x$, and $y_{1}$ satisfies both boundary
conditions. Since $P_{0}=1$ and
$$
W=
\left|\begin{array}{crccccc}
\sin x&\cos x \\\cos x&-\sin x
\end{array}\right|=-1,
$$
\eqref{eq:13.1.23} reduces to
$$
\int_{0}^{\pi} F(x)\sin x\,dx=0.
$$
From Example~\ref{example:13.1.2}, $F(x)=1$ and
$$
\int_{0}^{\pi} F(x)\sin x\,dx=\int_{0}^{\pi} \sin x\,dx=2,
$$
so      Theorem~\ref{thmtype:13.1.3}  implies that \eqref{eq:13.1.25} has no
solution.  In Example~\ref{example:13.1.3},
$$
F(x)=\sin 2x=2\sin x \cos x
$$
and
$$
\int_{0}^{\pi}F(x)\sin x\,dx=2\int_{0}^{\pi}\sin^{2}x\cos x\,dx
=\frac{2}{3}\sin^{3}x\bigg|_{0}^{\pi}=0,
$$
so Theorem~\ref{thmtype:13.1.3} implies that \eqref{eq:13.1.25} has infinitely
many solutions, differing by constant multiples of $y_{1}(x)=\sin x$.
\end{example}

\exercises

\begin{exerciselist}

\item\label{exer:13.1.1}Verify that $B_{1}$ and $B_{2}$ are linear operators; that is, if $c_{1}$
and $c_{2}$  are constants then
$$
B_{i}(c_{1}y_{1}+c_{2}y_{2})=c_{1}B_{i}(y_{1})+c_{2}B_{i}(y_{2}),\quad
i=1,2.
$$

\exercisetext{In Exercises~\ref{exer:13.1.2}--\ref{exer:13.1.7} solve the boundary
value problem.}

\item\label{exer:13.1.2}$y''-y=x$, \quad  $y(0)=-2$,\quad  $y(1)=1$


\item\label{exer:13.1.3}$y''=2-3x$, \quad  $y(0)=0$,\quad  $y(1)-y'(1)=0$

\item\label{exer:13.1.4}$y''-y=x$, \quad  $y(0)+y'(0)=3$,\quad  $y(1)-y'(1)=2$

\item\label{exer:13.1.5}$y''+4y=1$, \quad  $y(0)=3$,\quad
 $y(\pi/2)+y'(\pi/2)=-7$

\item\label{exer:13.1.6}$y''-2y'+y=2e^{x}$, \quad  $y(0)-2y'(0)=3$,\quad
 $y(1)+y'(1)=6e$

\item\label{exer:13.1.7}$y''-7y'+12y=4e^{2x}$, \quad  $y(0)+y'(0)=8$,\quad
 $y(1)=-7e^{2}$ (see Example~\ref{example:13.1.5})

\item\label{exer:13.1.8}State a condition on $F$ such that the boundary
value problem
$$
y''=F(x), \quad  y(0)=0, \quad y(1)-y'(1)=0
$$
has a  solution, and find all solutions.

\item\label{exer:13.1.9}\begin{alist}
\item % (a)
State  a condition on $a$ and $b$ such that the boundary
value problem
\begin{equation}
y''+y=F(x),\quad y(a)=0,\quad y(b)=0     \tag{A}
\end{equation}
has a unique solution for every continuous $F$, and find the solution
by the method used to prove Theorem~\ref{thmtype:13.1.3}
\item % (b)
In the case where $a$ and $b$ don't satisfy the condition you gave for
\part{a}, state necessary and sufficient     on $F$ such that (A) has
a solution, and find all solutions by the method used to prove
Theorem~\ref{thmtype:13.1.4}.
\end{alist}

\item\label{exer:13.1.10}Follow the instructions in Exercise~\ref{exer:13.1.9} for the boundary value
problem
$$
y''+y=F(x),\quad y(a)=0,\quad y'(b)=0.
$$

\item\label{exer:13.1.11}Follow the instructions in Exercise~\ref{exer:13.1.9} for the boundary value
problem
$$
y''+y=F(x),\quad y'(a)=0,\quad y'(b)=0.
$$

\exercisetext{In Exercises~\ref{exer:13.1.12}--\ref{exer:13.1.15} find a formula for
the solution of the  boundary problem by the method used to prove
Theorem~\ref{thmtype:13.1.3}. Assume that $a<b$.}

\item\label{exer:13.1.12}$y''-y=F(x)$,\quad  $y(a)=0$,\quad  $y(b)=0$

\item\label{exer:13.1.13}$y''-y=F(x)$,\quad  $y(a)=0$,\quad  $y'(b)=0$

\item\label{exer:13.1.14}$y''-y=F(x)$,\quad  $y'(a)=0$,\quad  $y'(b)=0$

\item\label{exer:13.1.15}$y''-y=F(x)$,\quad  $y(a)-y'(a)=0$,\quad  $y(b)+y'(b)=0$

\exercisetext{In Exercises~\ref{exer:13.1.16}--\ref{exer:13.1.19} find all values of
$\omega$ such that
 boundary problem has a unique solution, and find the solution  by the
method used to prove Theorem~\ref{thmtype:13.1.3}. For other values of $\omega$,
find conditions on $F$ such that the problem has a solution,
 and find all solutions  by the
method used to prove Theorem~\ref{thmtype:13.1.4}.}

\item\label{exer:13.1.16}$y''+ \omega^{2}y=F(x)$,\quad $y(0)=0$,\quad $y(\pi)=0$

\item\label{exer:13.1.17}$y''+ \omega^{2}y=F(x)$,\quad $y(0)=0$,\quad $y'(\pi)=0$

\item\label{exer:13.1.18}$y''+ \omega^{2}y=F(x)$,\quad $y'(0)=0$,\quad $y(\pi)=0$

\item\label{exer:13.1.19}$y''+ \omega^{2}y=F(x)$,\quad $y'(0)=0$,\quad $y'(\pi)=0$

\item\label{exer:13.1.20}Let $\{z_{1},z_{2}\}$ be a fundamental set of solutions of $Ly=0$.
Given that the homogeneous boundary value problem
$$
Ly=0, \quad B_{1}(y)=0,\quad B_{2}(y)=0
$$
has a nontrivial solution, express it explicity in terms of $z_{1}$
and $z_{2}$.

\item\label{exer:13.1.21}If the boundary value problem has a solution for every continuous  $F$,
then find  the Green's function for the problem and use it to write
an explicit formula for the solution. Otherwise, if the
boundary value problem does not have a solution for every continuous $F$,
find a necessary and sufficient condition on $F$ for the problem to have
a solution, and find all solutions. Assume that $a<b$.

\begin{alist}
\item % (a)
$y''=F(x)$, \quad $y(a)=0$,\quad $y(b)=0$

\item % (b)
$y''=F(x)$, \quad $y(a)=0$,\quad $y'(b)=0$

\item % (c)
$y''=F(x)$, \quad $y'(a)=0$,\quad $y(b)=0$

\item % (d)
$y''=F(x)$, \quad $y'(a)=0$,\quad $y'(b)=0$
\end{alist}

\item\label{exer:13.1.22}Find the Green's function  for the boundary value problem
\begin{equation}
y''=F(x), \quad y(0)-2y'(0)=0, \quad y(1)+2y'(1)=0.   \tag{A}
\end{equation}
Then  use the Green's function to solve (A) with
\part{a} $F(x)=1$, \part{b} $F(x)=x$, and \part{c} $F(x)=x^{2}$.

\item\label{exer:13.1.23}Find the Green's function  for the boundary value problem
\begin{equation}
x^{2}y''+xy'+(x^{2}-1/4)y=F(x), \quad y(\pi/2)=0,\quad y(\pi)=0,   \tag{A}
\end{equation}
given that
$$
y_{1}(x)=\frac{\cos x}{\sqrt{x}}
\text{\;    and\; \;  }
y_{2}(x)=\frac{\sin x}{\sqrt{x}}
$$
are solutions of the complementary equation. Then use the Green's function
to solve  (A) with \part{a} $F(x)=x^{3/2}$  and \part{b} $F(x)=x^{5/2}$.

\item\label{exer:13.1.24}Find the Green's function  for the boundary value problem
\begin{equation}
x^{2}y''-2xy'+2y=F(x), \quad y(1)=0,\quad y(2)=0,   \tag{A}
\end{equation}
given that  $\{x,x^{2}\}$ is a fundamental set of solutions
of the complementary equation.
Then use the Green's function
to solve  (A) with \part{a} $F(x)=2x^{3}$  and \part{b} $F(x)~=~6x^{4}$.


\item\label{exer:13.1.25}Find the Green's function  for the boundary value problem
\begin{equation}
x^{2}y''+xy'-y=F(x), \quad y(1)-2y'(1)=0,\quad y'(2)=0,   \tag{A}
\end{equation}
given that  $\{x,1/x\}$ is a fundamental set of solutions
of the complementary equation.
Then use the Green's function
to solve  (A) with \part{a} $F(x)=1$,   \part{b} $F(x)=x^{2}$, and
\part{c} $F(x)=x^{3}$.

\exercisetext{In Exercises~\ref{exer:13.1.26}--\ref{exer:13.1.30} find
necessary and sufficient conditions on $\alpha$, $\beta$, $\rho$,
and $\delta$ for the boundary value
 problem to have a  unique solution for every continuous $F$, and
find the Green's function.}

\item\label{exer:13.1.26}$y''=F(x)$, \quad $\alpha y(0)+\beta y'(0)=0$, \quad
$\rho y(1)+\delta y'(1)=0$

\item\label{exer:13.1.27}$y''+y=F(x)$, \quad $\alpha y(0)+\beta y'(0)=0$, \quad
$\rho y(\pi)+\delta y'(\pi)=0$

\item\label{exer:13.1.28}$y''+y=F(x)$, \quad $\alpha y(0)+\beta y'(0)=0$, \quad
$\rho y(\pi/2)+\delta y'(\pi/2)=0$

\item\label{exer:13.1.29}$y''-2y'+2y=F(x)$, \quad $\alpha y(0)+\beta y'(0)=0$, \quad
$\rho y(\pi)+\delta y'(\pi)=0$

\item\label{exer:13.1.30}$y''-2y'+2y=F(x)$, \quad $\alpha y(0)+\beta y'(0)=0$, \quad
$\rho y(\pi/2)+\delta y'(\pi/2)=0$

\item\label{exer:13.1.31}Find
necessary and sufficient conditions on $\alpha$, $\beta$, $\rho$,
and $\delta$ for the boundary value
 problem
\begin{equation}
y''-y=F(x), \quad \alpha y(a)+\beta y'(a)=0, \quad
\rho y(b)+\delta y'(b)=0   \tag{A}
\end{equation}
to have a  unique solution for every continuous $F$, and find the Green's
function for (A). Assume that $a<b$.

\item\label{exer:13.1.32}Show that  the assumptions of Theorem~\ref{thmtype:13.1.3} imply that the unique
solution of
$$
Ly=F, \quad  B_{1}(y)=k_{1},\quad B_{2}(y)=f_{2}
$$
is
$$
y=\int_{a}^{b} G(x,t)F(t)\,dt +\frac{k_{2}}{B_{2}}(y_{1})y_{1}
+\frac{k_{1}}{B_{1}(y_{2})}y_{2}.
$$
\end{exerciselist}


\newsection{2}
{Boundary Value Problems for Second Order Ordinary
Differential Equations}
{Sturm-Liouville Problems}
\currentpdfbookmark{Section 13.2 Sturm-Liouville Problems} {section:13.2}
\vskip14pt


\renewcommand{\thissection}
{\sectiontitle{STURM-LIOUVILLE PROBLEMS}}
\thissection


\noindent
In this section we consider eigenvalue problems of  the form
\begin{equation} \label{eq:13.2.1}
P_{0}(x)y''+P_{1}(x)y'+P_{2}(x)y+ \lambda R(x)y=0,\quad B_{1}(y)=0,\quad
B_{2}(y)=0,
\end{equation}
where
$$
B_{1}(y)=\alpha y(a)+\beta y'(a) \text{\quad and\quad}
B_{2}(y)=\rho y(b)+\delta y'(b).
$$
As in Section~13.1,  $\alpha$, $\beta$, $\rho$, and $\delta$
are real numbers, with
$$
\alpha^{2}+\beta^{2}>0\text{\;    and\; \;  } \rho^{2}+\delta^{2}>0,
$$
$P_{0}$, $P_{1}$, $P_{2}$, and $R$ are continuous, and $P_{0}$ and $R$
are positive on $[a,b]$.





We say that $\lambda$ is an \emph{eigenvalue} of  \eqref{eq:13.2.1} if
\eqref{eq:13.2.1} has a nontrivial solution $y$.   In this case, $y$ is an
\emph{eigenfunction associated with $\lambda$}, or
a $\lambda$-\emph{eigenfunction}.   \emph{Solving}  the eigenvalue
problem means finding all eigenvalues and associated eigenfunctions
of \eqref{eq:13.2.1}.


\begin{example}\label{example:13.2.1} \rm
Solve the eigenvalue problem
\begin{equation} \label{eq:13.2.2}
y'+3y'+2y+\lambda y=0,\quad  y(0)=0,\quad y(1)=0.
\end{equation}
\end{example}

\solution
The characteristic equation of \eqref{eq:13.2.2} is
$$
r^{2}+3r+2+\lambda=0,
$$
with zeros
$$
r_{1}=\frac{-3+\sqrt{1-4\lambda}}{2} \text{\quad   and \quad }
r_{2}=\frac{-3-\sqrt{1-4\lambda}}{2}.
$$
If $\lambda<1/4$ then  $r_{1}$ and $r_{2}$ are real and distinct, so the
general solution of the differential equation in  \eqref{eq:13.2.2} is
$$
y=c_{1}e^{r_{1}t}+c_{2}e^{r_{2}t}.
$$
The boundary conditions require that
\begin{eqnarray*}
c_{1}\phantom{e^{r_{1}}}+c_{2}\phantom{e^{r_{2}}}&=&0\\
c_{1}e^{r_{1}}+c_{2}e^{r_{2}}&=&0.
\end{eqnarray*}
Since the determinant of  this system is $e^{r_{2}}-e^{r_{1}}\ne0$, the
system has only the trivial solution. Therefore
$\lambda$ isn't  an eigenvalue of \eqref{eq:13.2.2}.

If $\lambda=1/4$ then $r_{1}=r_{2}=-3/2$, so the general solution of
\eqref{eq:13.2.2} is
$$
y=e^{-3x/2}(c_{1}+c_{2}x).
$$
The boundary condition $y(0)=0$ requires that $c_{1}=0$, so
$y=c_{2}xe^{-3x/2}$ and the boundary condition $y(0)$ requires that
$c_{2}=0$. Therefore $\lambda=1/4$  isn't  an eigenvalue of~\eqref{eq:13.2.2}.

If $\lambda>1/4$ then
$$
r_{1}=-\frac{3}{2}+i\omega \text{\quad   and\quad }
r_{2}=-\frac{3}{2}-i\omega,
$$
with
\begin{equation} \label{eq:13.2.3}
\omega =\frac{\sqrt{4\lambda-1}}{2}
\text{\quad  or, equivalently,\quad }  \lambda=\frac{1+4\omega^{2}}{4}.
\end{equation}
In this case the general solution of the differential equation in
\eqref{eq:13.2.2} is
$$
y=e^{-3x/2}(c_{1}\cos\omega x+c_{2}\sin\omega x).
$$
The boundary condition $y(0)=0$  requires that $c_{1}=0$, so
$y=c_{2}e^{-3x/2}\sin\omega x$, which holds with $c_{2}\ne0$ if and only if
$\omega=n\pi$, where $n$ is an integer. We may assume that $n$
is a positive integer. (Why?). From \eqref{eq:13.2.3}, the eigenvalues are
$\lambda_{n}=(1+4n^{2}\pi^{2})/4$, with associated eigenfunctions
$$
y_{n}=e^{-3x/2}\sin n\pi x,\quad  n=1,2,3,\dots.
$$


\begin{example} \label{example:13.2.2} \rm
Solve the eigenvalue problem
\begin{equation} \label{eq:13.2.4}
x^{2}y''+xy'+\lambda y=0,\quad y(1)=0,\quad y(2)=0.
\end{equation}
\end{example}

\solution
If  $\lambda=0$,  the differential equation in \eqref{eq:13.2.4}  reduces to
$x(xy')'=0$, so $xy'=c_{1}$,
$$
y'=\frac{c_{1}}{x}, \text{\quad and\quad}  y=c_{1}\ln x+c_{2}.
$$
The boundary condition $y(1)=0$ requires that $c_{2}=0$, so $y=c_{1}\ln x$.
The boundary condition  $y(2)=0$ requires that $c_{1}\ln2=0$, so $c_{1}=0$.
Therefore zero isn't  an eigenvalue of \eqref{eq:13.2.4}.

 If $\lambda<0$, we
write $\lambda=-k^{2}$ with $k>0$, so \eqref{eq:13.2.4} becomes
$$
x^{2}y''+xy'-k^{2}y=0,
$$
an Euler equation (Section~7.4) with indicial equation
$$
r^{2}-k^{2}=(r-k)(r+k)=0.
$$
Therefore
$$
y=c_{1}x^{k}+c_{2}x^{-k}.
$$
The boundary conditions require that
\begin{eqnarray*}
\phantom{2^{k}}c_{1}+\phantom{2^{-k}}c_{2}&=&0 \\
2^{k}c_{1}+2^{-k}c_{2}&=&0.
\end{eqnarray*}
Since the determinant of this system  is $2^{-k}-2^{k}\ne0$,
$c_{1}=c_{2}=0$. Therefore \eqref{eq:13.2.4} has no negative eigenvalues.

If $\lambda>0$ we write $\lambda=k^{2}$ with $k>0$. Then  \eqref{eq:13.2.4}
becomes
$$
x^{2}y''+xy' +k^{2}y=0,
$$
an Euler equation with indicial equation
$$
r^{2}+k^{2}=(r-ik)(r+ik)=0,
$$
so
$$
y=c_{1}\cos(k\ln x)+c_{2}\sin(k\ln x).
$$
 The boundary condition $y(1)=0$
requires that $c_{1}=0$. Therefore $y=c_{2}\sin(k\ln x)$. This holds with
$c_{2}\ne0$ if and only if $k=n\pi/\ln 2$, where $n$ is a positive integer.
Hence, the eigenvalues of \eqref{eq:13.2.4} are $\lambda_{n}=(n\pi/\ln2)^{2}$,
with associated eigenfunctions
$$
y_{n}=\sin\left(\frac{n\pi}{\ln2}\ln x\right),\quad n=1,2,3,\dots.
$$

For theoretical purposes, it's useful to rewrite the differential equation
in \eqref{eq:13.2.1} in a different form, provided by the next theorem.

\begin{theorem}\color{blue} \label{thmtype:13.2.1}
If $P_{0},$ $P_{1},$ $P_{2},$ and $R$ are continuous and $P_{0}$ and $R$
are positive on a closed interval $[a,b],$ then the equation
\begin{equation} \label{eq:13.2.5}
P_{0}(x)y''+P_{1}(x)y'+P_{2}(x)y+\lambda R(x)y=0
\end{equation}
can be rewritten as
\begin{equation} \label{eq:13.2.6}
(p(x)y')'+q(x)y+\lambda r(x)y=0,
\end{equation}
where $p$, $p'$, $q$ and $r$ are continuous and $p$ and $r$ are positive
on $[a,b].$
\end{theorem}

\proof
We begin by rewriting \eqref{eq:13.2.5} as
\begin{equation} \label{eq:13.2.7}
y''+u(x)y'+v(x)y+\lambda R_{1}(x)y=0,
\end{equation}
with $u=P_{1}/P_{0}$, $v=P_{2}/P_{0}$, and $R_{1}=R/P_{0}$. (Note that
$R_{1}$ is positive on $[a,b]$.)  Now let $p(x)=e^{U(x)}$, where $U$ is any
antiderivative of $u$. Then $p$ is positive  on  $[a,b]$ and, since
$U'=u$,
\begin{equation} \label{eq:13.2.8}
p'(x)=p(x)u(x)
\end{equation}
is continuous on $[a,b]$. Multiplying \eqref{eq:13.2.7} by $p(x)$  yields
\begin{equation} \label{eq:13.2.9}
p(x)y''+p(x)u(x)y'+p(x)v(x)y+\lambda p(x)R_{1}(x)y=0.
\end{equation}
Since $p$ is positive on $[a,b]$, this equation has the same solutions as
\eqref{eq:13.2.5}. From \eqref{eq:13.2.8},
$$
(p(x)y')'=p(x)y''+p'(x)y'=p(x)y''+p(x)u(x)y',
$$
so \eqref{eq:13.2.9} can be rewritten as in \eqref{eq:13.2.6}, with $q(x)=p(x)v(x)$
and $r(x)=p(x)R_{1}(x)$. This completes the proof.

It is to be understood throughout the rest of this section that $p$, $q$,
and $r$ have the properties stated in Theorem~\ref{thmtype:13.2.1}. Moreover,
whenever we write $Ly$ in a general statement, we mean
$$
Ly=(p(x)y')'+q(x)y.
$$

The differential equation \eqref{eq:13.2.6} is called a
\href{http://www-history.mcs.st-and.ac.uk/Mathematicians/Sturm.html}
{\color{blue}\it Sturm}--\href{http://www-history.mcs.st-and.ac.uk/Mathematicians/Liouville.html}
{\color{blue}\it Liouville  equation}, and the
eigenvalue problem
\begin{equation} \label{eq:13.2.10}
(p(x)y')'+q(x)y+\lambda r(x)y=0,\quad B_{1}(y)=0,\quad
B_{2}(y)=0,
\end{equation}
which is equivalent to \eqref{eq:13.2.1}, is called a
\href{http://www-history.mcs.st-and.ac.uk/Mathematicians/Sturm.html}
{\color{blue}\it Sturm-Liouville
problem}.

\begin{example}\label{example:13.2.3} \rm
Rewrite the eigenvalue problem
\begin{equation} \label{eq:13.2.11}
y''+3y'+(2+\lambda)y=0,\quad y(0)=0,\quad y(1)=0
\end{equation}
of Example~\ref{example:13.2.1}  as a Sturm-Liouville problem.
\end{example}

\solution
Comparing \eqref{eq:13.2.11} to \eqref{eq:13.2.7}  shows that $u(x)=3$, so we   take
$U(x)=3x$ and $p(x)=e^{3x}$. Multiplying the differential equation in
\eqref{eq:13.2.11} by $e^{3x}$ yields
$$
e^{3x}(y''+3y')+2e^{3x}y+\lambda e^{3x}y=0.
$$
Since
$$
e^{3x}(y''+3y')=(e^{3x}y')',
$$
 \eqref{eq:13.2.11} is equivalent to the Sturm--Liouville problem
\begin{equation} \label{eq:13.2.12}
(e^{3x}y')'+2e^{3x}y+\lambda e^{3x}y=0,\quad y(0)=0,\quad y(1)=0.
\end{equation}


\begin{example}\label{example:13.2.4} \rm
Rewrite the eigenvalue problem
\begin{equation} \label{eq:13.2.13}
x^{2}y''+xy'+\lambda y=0,\quad y(1)=0,\quad y(2)=0
\end{equation}
of Example~\ref{example:13.2.2}  as a Sturm-Liouville problem.
\end{example}



\solution
Dividing the differential equation in \eqref{eq:13.2.13} by $x^{2}$  yields
$$
y''+\frac{1}{x}y'+\frac{\lambda}{x^{2}}y=0.
$$
Comparing this to \eqref{eq:13.2.7} shows that $u(x)=1/x$, so we take $U(x)=\ln
x$ and $p(x)=e^{\ln x}=x$.
Multiplying the differntial equation by $x$ yields
$$
xy''+y'+\frac{\lambda}{x}y=0.
$$
Since
$$
xy''+y'=(xy')',
$$
\eqref{eq:13.2.13} is equivalent to the Sturm--Liouville problem
\begin{equation} \label{eq:13.2.14}
(xy')'+\frac{\lambda}{x}y=0,\quad y(1)=0,\quad y(2)=0.
\end{equation} \bbox

Problems 1--4 of Section~11.1  are Sturm--Liouville problems.
(Problem~5 isn't , although some authors use a definition of
\href{http://www-history.mcs.st-and.ac.uk/Mathematicians/Sturm.html}
{\color{blue}\it Sturm-Liouville problem}
that does include it.)  We were able to find
the eigenvalues of Problems 1-4  explicitly because in each problem the
coefficients in the boundary conditions satisfy $\alpha\beta=0$ and
$\rho\delta=0$; that is, each boundary condition  involves either $y$
or $y'$, but not both. If this isn't true then the eigenvalues can't in
general be expressed exactly by simple formulas;     rather, approximate
values
must be obtained by numerical solution of equations derived by requiring
the determinants of certain $2\times 2$ systems of homogeneous equations to
be zero. To apply the numerical methods effectively, graphical methods
must be used to determine approximate locations of the zeros of these
determinants. Then the zeros can be computed accurately by
numerical methods.

\begin{example}  \label{example:13.2.5} \rm
Solve the Sturm--Liouville problem
\begin{equation} \label{eq:13.2.15}
y''+\lambda y=0, \quad y(0)+y'(0)=0,\quad y(1)+3y'(1)=0.
\end{equation}
\end{example}

\solution  If
$\lambda=0$, the differential equation in \eqref{eq:13.2.15} reduces
to $y''=0$, with general solution $y=c_{1}+c_{2}x$. The boundary conditions
require that
\begin{eqnarray*}
c_{1}+\phantom{4}c_{2}&=&0\\
c_{1}+4c_{2}&=&0,
\end{eqnarray*}
so $c_{1}=c_{2}=0$. Therefore zero isn't  an eigenvalue of \eqref{eq:13.2.15}.

If $\lambda<0$,  we write
$\lambda=-k^{2}$ where $k>0$, and the differential equation in
\eqref{eq:13.2.15}   becomes $y''-k^{2}y=0$, with general solution
\begin{equation} \label{eq:13.2.16}
y=c_{1}\cosh kx+c_{2}\sinh kx,
\end{equation}
so
$$
y'=k(c_{1}\sinh kx+c_{2}\cosh kx).
$$
The boundary conditions require that
\begin{equation} \label{eq:13.2.17}
\begin{gathered}
c_{1}+kc_{2}=0\\
(\cosh k+3k\sinh k)c_{1}+(\sinh k+3k\cosh k)c_{2}=0.
\end{gathered}
\end{equation}
The determinant of this system is
\begin{eqnarray*}
D_{N}(k)&=&
\left|\begin{array}{cccccc}
1&k\\
\cosh k+3k\sinh k&\sinh k+3k \cosh k
\end{array}\right|      \\
&=& (1-3k^{2})\sinh k+2k \cosh k.
\end{eqnarray*}
Therefore the  system \eqref{eq:13.2.17} has a nontrivial solution if and only
if
$D_{N}(k)=0$ or, equivalently,
\begin{equation} \label{eq:13.2.18}
\tanh k=-\frac{2k}{1-3k^{2}}.
\end{equation}
The graph of the right side (Figure~\ref{figure:13.2.1}) has a vertical
asymptote at
$k=1/\sqrt{3}$.
Since the two sides have different signs if $k<1/\sqrt{3}$,  this equation
has no solution in $(0,1/\sqrt{3})$. Figure~\ref{figure:13.2.1} shows the
graphs
of the two sides of \eqref{eq:13.2.18} on an interval to the right of the
vertical asymptote, which is indicated by the dashed line. You can see that
the two curves intersect near $k_{0}=1.2$, Given this estmate, you can use
Newton's to compute $k_{0}$  more accurately. We computed $k_{0}\approx
1.1219395$. Therefore $-k_{0}^{2}\approx -1.2587483$ is an eigenvalue of
\eqref{eq:13.2.15}. From \eqref{eq:13.2.16} and the first equation in \eqref{eq:13.2.17},
$$
y_{0}=k_{0}\cosh k_{0}x-\sinh k_{0}x.
$$


\begin{figure}[tbp]
  \centering\color{blue}
  \includegraphics[bb=-78 148 689 643,width=5.67in,height=3.66in,keepaspectratio]{fig130201}
\caption{$u=\tanh k$  and $u=-2k/(1-3k^2)$}
  \label{figure:13.2.1}
\end{figure}




If $\lambda>0$ we write
$\lambda=k^{2}$ where $k>0$, and   differential equation in
\eqref{eq:13.2.15} becomes $y''+k^{2}y=0$, with general solution
\begin{equation} \label{eq:13.2.19}
y= \cos kx + c_{2}\sin kx,
\end{equation}
so
$$
y'=k(-c_{1}\sin kx+c_{2}\cos kx).
$$
The boundary conditions require that
\begin{equation} \label{eq:13.2.20}
\begin{gathered}
c_{1}+kc_{2}=0\\
(\cos k-3k\sin k)c_{1}+(\sin k+3k\cos k)c_{2}=0.
\end{gathered}
\end{equation}
The determinant of this system is
\begin{eqnarray*}
D_{P}(k)&=&
\left|\begin{array}{cccccc}
1&k\\
\cos k-3k\sin k&\sin k+3k \cos k
\end{array}\right|      \\
&=& (1+3k^{2})\sin k+2k \cos k.
\end{eqnarray*}
The system \eqref{eq:13.2.20}  has a nontrivial solution if and only if
$D_{P}(k)=0$ or, equivalently,
$$
\tan k=-\frac{2k}{1+3k^{2}}.
$$
Figure~\ref{figure:13.2.2} shows  the graphs of the two sides of this
equation.
You can see from the figure that  the graphs intersect at infinitely many
points
$k_{n}\approx n\pi$ ($n=1$, $2$, $3$,\dots), where the error in this
approximation approaches zero as $n\to\infty$. Given this estimate, you can
use Newton's method to compute $k_{n}$  more accurately. We computed
\begin{eqnarray*}
k_{1}&\approx&\phantom{1}2.9256856,\\
k_{2}&\approx&\phantom{1} 6.1765914,\\
k_{3}&\approx&\phantom{1} 9.3538959,\\
k_{4}&\approx& 12.5132570.
\end{eqnarray*}
The estimates of the corresponding eigenvalues $\lambda_{n}=k_{n}^{2}$ are
\begin{eqnarray*}
\lambda_{1}&\approx&\phantom{15} 8.5596361,\\
\lambda_{2}&\approx&\phantom{5} 38.1502809,\\
\lambda_{3}&\approx&\phantom{5} 87.4953676,\\
\lambda_{4}&\approx& 156.5815998.
\end{eqnarray*}
From \eqref{eq:13.2.19}   and the first equation in \eqref{eq:13.2.20},
$$
y_{n}=k_{n}\cos k_{n}x-\sin k_{n}x
$$
is an eigenfunction associated with $\lambda_{n}$  \bbox



\begin{figure}[tbp]
  \centering\color{blue}
  \includegraphics[bb=-78 148 689 643,width=5.67in,height=3.66in,keepaspectratio]{fig130202}
\caption{$u=\tan k$ and $u=-2k/(1+k)$}
  \label{figure:13.2.2}
\end{figure}




Since the differential equations in \eqref{eq:13.2.12} and \eqref{eq:13.2.14}
are more complicated than those in \eqref{eq:13.2.11} and \eqref{eq:13.2.13}
respectively, what is the point of Theorem~\ref{thmtype:13.2.1}?
The point is this: to solve a \emph{specific} problem, it may be better
to deal with it directly, as we did in Examples~\ref{example:13.2.1}
and \ref{example:13.2.2};     however, we'll see that transforming the
\emph{general} eigenvalue problem \eqref{eq:13.2.1} to the Sturm--Liouville
problem \eqref{eq:13.2.10}  leads to results applicable to \emph{all}
eigenvalue problems of the form \eqref{eq:13.2.1}.

\begin{theorem}\color{blue} \label{thmtype:13.2.2}
If
$$
Ly=(p(x)y')'+q(x)y
$$
and $u$ and $v$  are twice continuously  functions on $[a,b]$ that satisfy
the boundary conditions $B_{1}(y)=0$ and $B_{2}(y)=0,$  then
\begin{equation} \label{eq:13.2.21}
\int_{a}^{b}[u(x)Lv(x)-v(x)Lu(x)]\,dx=0.
\end{equation}
\end{theorem}

\proof
Integration by parts yields
\begin{eqnarray*}
\int_{a}^{b}[u(x)Lv(x)-v(x)Lu(x)]\,dx&=&
\int_{a}^{b}[u(x)(p(x)v'(x))'-v(x)(p(x)u'(x))']\,dx\\
&=& p(x)[u(x)v'(x)-u'(x)v(x)]\bigg|_{a}^{b}\\
&&-\int_{a}^{b}p(x)[u'(x)v'(x)-u'(x)v'(x)]\,dx.
\end{eqnarray*}
Since the last integral equals zero,
\begin{equation} \label{eq:13.2.22}
\int_{a}^{b}[u(x)Lv(x)-v(x)Lu(x)]\,dx =
 p(x)[u(x)v'(x)-u'(x)v'(x)]\bigg|_{a}^{b}.
\end{equation}
By assumption, $B_{1}(u)=B_{1}(v)=0$ and $B_{2}(u)=B_{2}(v)=0$. Therefore
$$
\begin{aligned}
\alpha u(a)+\beta u'(a)&=0\\
\alpha v(a)+\beta v'(a)&=0\\
\end{aligned}
\quad \text{\quad   and\quad}\quad
\begin{gathered}
\rho u(b)+\delta u'(b)=0\phantom{.}\\
\rho v(b)+\delta v'(b)=0.
\end{gathered}
$$
Since $\alpha^{2}+\beta^{2}>0$ and $\rho^{2}+\delta^{2}>0$, the
determinants of these two systems must both be zero; that is,
$$
u(a)v'(a)-u'(a)v(a)=u(b)v'(b)-u'(b)v(b)=0.
$$
This and \eqref{eq:13.2.22} imply \eqref{eq:13.2.21}, which completes the
proof.

The next theorem shows that a Sturm--Liouville problem has no complex
eigenvalues.

\begin{theorem}\color{blue} \label{thmtype:13.2.3}
If $\lambda=p+qi$ with $q\ne0$  then the boundary value problem
$$
Ly+\lambda r(x)y=0,\quad B_{1}(y)=0,\quad B_{2}(y)=0
$$
has only the trivial solution.
\end{theorem}

\proof For this theorem to make sense, we must consider complex-valued
solutions  of
\begin{equation} \label{eq:13.2.23}
Ly+(p+iq)r(x,y)y=0.
\end{equation}
If $y=u+iv$ where $u$ and $v$ are real-valued and twice differentiable,
we define $y'=u'+iv'$ and $y''=u''+iv''$. We say that $y$  is a solution of
\eqref{eq:13.2.23} if
the   real and imaginary parts of the left side of \eqref{eq:13.2.23} are both
zero.  Since $Ly=(p(x)'y)'+q(x)y$  and $p$, $q$, and $r$ are real-valued,
\begin{eqnarray*}
Ly+\lambda
r(x)y&=&L(u+iv)+(p+iq)r(x)(u+iv)\\
&=&Lu+r(x)(pu-qv)+i[Lv+r(x)(pu+qv)],
\end{eqnarray*}
so $Ly+\lambda r(x)y=0$ if and only if
\begin{eqnarray*}
Lu+r(x)(pu-qv)&=&0\\
Lv+r(x)(qu+pv)&=&0.
\end{eqnarray*}
Multiplying the first equation by $v$ and the second by $u$  yields
\begin{eqnarray*}
vLu+r(x)(puv-qv^{2})&=&0\\
uLv+r(x)(qu^{2}+puv)&=&0.
\end{eqnarray*}
Subtracting the first equation from the second yields
$$
uLv-vLu+qr(x)(u^{2}+v^{2})=0,
$$
so
\begin{equation} \label{eq:13.2.24}
\int_{a}^{b}[u(x)Lv(x)-v(x)Lu(x)]\,dx+
\int_{a}^{b}r(x)[u^{2}(x)+v^{2}(x)]\,dx=0.
\end{equation}
Since
$$
B_{1}(y)=B_{1}(u+iv)=B_{1}(u)+iB_{1}(v)
$$
and
$$
B_{2}(y)=B_{2}(u+iv)=B_{2}(u)+iB_{2}(v),
$$
$B_{1}(y)=0$ and $B_{2}(y)=0$  implies that
$$
B_{1}(u)=B_{2}(u)=B_{1}(v)=B_{2}(v)=0.
$$
 Therefore Theorem~\ref{thmtype:13.2.2}
implies that  first integral in
\eqref{eq:13.2.24} equals zero, so \eqref{eq:13.2.24} reduces to
$$
q\int_{a}^{b}r(x)[u^{2}(x)+v^{2}(x)]\,dx =0.
$$
Since $r$ is positive on $[a,b]$  and $q\ne0$ by assumption, this implies
that $u\equiv0$ and $v\equiv0$ on $[a,b]$. Therefore $y\equiv0$ on
$[a,b]$, which completes the proof.

\begin{theorem}\color{blue} \label{thmtype:13.2.4}
If $\lambda_{1}$ and $\lambda_{2}$ are distinct eigenvalues of the
Sturm--Liouville  problem
\begin{equation} \label{eq:13.2.25}
Ly+\lambda r(x)y=0,\quad B_{1}(y)=0,\quad B_{2}(y)=0
\end{equation}
with associated eigenfunctions $u$ and $v$ respectively$,$ then
\begin{equation} \label{eq:13.2.26}
\int_{a}^{b}r(x)u(x)v(x)\,dx=0.
\end{equation}
\end{theorem}


\proof Since $u$ and $v$ satisfy the boundary conditions in \eqref{eq:13.2.25},
Theorem~\ref{thmtype:13.2.2} implies that
$$
\int_{a}^{b}[u(x)Lv(x)-v(x)Lu(x)]\,dx=0.
$$
Since $Lu=-\lambda_{1}ru$ and $Lv=-\lambda_{2}rv$, this implies that
$$
(\lambda_{1}-\lambda_{2})\int_{a}^{b}r(x)u(x)v(x)\,dx=0.
$$
Since $\lambda_{1}\ne\lambda_{2}$, this implies \eqref{eq:13.2.26}, which
completes the proof.

If $u$  and $v$ are any integrable functions on $[a,b]$ and
$$
\int_{a}^{b} r(x)u(x)v(x)\,dx=0,
$$
we say that $u$ and $v$ \emph{orthogonal on} $[a,b]$
 \emph{with respect to} $r=r(x)$.

Theorem~\ref{thmtype:13.1.1}  implies the next theorem.

\begin{theorem}\color{blue} \label{thmtype:13.2.5}
If $u\not\equiv0$ and $v$ both satisfy
$$
Ly+\lambda r(x)y=0,\quad B_{1}(y)=0,\quad B_{2}(y)=0,
$$
then $v=cu$ for some constant $c.$
\end{theorem}


We've now proved parts of the next theorem. A complete proof is
beyond the scope of this book.


\begin{theorem}\color{blue} \label{thmtype:13.2.6}
The set of all eigenvalues of the Sturm--Liouville problem
$$
Ly+\lambda r(x)y=0,\quad B_{1}(y)=0,\quad B_{2}(y)=0
$$
can be ordered as
$$
\lambda_{1}<\lambda_{2}<\cdots<\lambda_{n}<\cdots,
$$
and
$$
\lim_{n\to\infty} \lambda_{n}=\infty.
$$
For each $n,$  if $y_{n}$ is an arbitrary $\lambda_{n}$-eigenfunction$,$
then every $\lambda_{n}$-eigenfunction is a constant multiple of $y_{n}.$
If $m\ne n,$   $y_{m}$ and $y_{n}$  are orthogonal $[a,b]$  with
respect to  $r=r(x);$    that is$,$
\begin{equation} \label{eq:13.2.27}
\int_{a}^{b} r(x)y_{m}(x)y_{n}(x)\,dx=0.
\end{equation}
\end{theorem}

 You may want to verify \eqref{eq:13.2.27} for the eigenfunctions obtained in
Examples~\ref{example:13.2.1}  and \ref{example:13.2.2}.


In conclusion, we mention the next theorem. The proof is beyond the
scope of this book.

\begin{theorem}\color{blue} \label{thmtype:13.2.7}
Let $\lambda_{1}<\lambda_{2}<\cdots<\lambda_{n}<\cdots$  be the eigenvalues
of the Sturm--Liouville problem
$$
Ly+\lambda r(x)y=0,\quad B_{1}(y)=0,\quad B_{2}(y)=0,
$$
with associated eigenvectors $y_{1},$  $y_{2},$ \dots, $y_{n},$ \dots$.$
 Suppose $f$  is piecewise smooth
$($Definition~$\ref{thmtype:11.2.3})$ on $[a,b].$  For each $n,$ let
$$
c_{n}=\frac{\dst{\int_{a}^{b}}r(x)f(x)y_{n}(x)\,dx}
{\dst{\int_{a}^{b}}r(x)y_{n}^{2}(x)\,dx}.
$$
Then
$$
\frac{f(x-)+f(x+)}{2}=\sum_{n=1}^{\infty}c_{n}y_{n}(x)
$$
for all $x$ in the open interval $(a,b).$
\end{theorem}

\exercises



\emph{In Exercises~$\ref{exer:13.2.1}$--$\ref{exer:13.2.7}$  rewrite the equation
in Sturm--Liouville form $($with $\lambda=0)$. Assume that $b,$ $c,$
$\alpha,$ and $\nu$ are contants$.$}

\begin{exerciselist}

\item\label{exer:13.2.1}$y''+by'+cy=0$

\item\label{exer:13.2.2}$x^{2}y''+xy'+(x^{2}-\nu^{2})y=0$ \quad (Bessel's
equation)

\item\label{exer:13.2.3}$(1-x^{2})y''-xy'+\alpha^{2}y=0$ \quad
(Chebyshev's equation)

\item\label{exer:13.2.4}$x^{2}y''+bxy'+cy=0$ \quad (Euler's equation)

\item\label{exer:13.2.5}$y''-2xy'+2\alpha y=0$ \quad (Hermite's equation)
\item\label{exer:13.2.6}$xy''+(1-x)y'+\alpha y=0$ \quad (Laguerre's equation)



\item\label{exer:13.2.7}$(1-x^{2})y''-2xy'+\alpha(\alpha+1)y=0$
\quad (Legendre's equation)




\item\label{exer:13.2.8}In Example~\ref{example:13.2.4} we found that the eigenvalue problem
\begin{equation}
x^{2}y''+xy'+\lambda y=0,\quad y(1)=0,\quad y(2)=0 \tag{A}
\end{equation}
is equivalent to the Sturm-Liouville problem
\begin{equation} \tag{B}
(xy')'+\frac{\lambda}{x}y=0,\quad y(1)=0,\quad y(2)=0.
\end{equation}
Multiply the differential equation in (B) by $y$ and integrate to show that
$$
\lambda\int_{1}^{2}\frac{y^{2}(x)}{x}\,dx=\int_{1}^{2}x(y'(x))^{2}\,dx.
$$
Conclude from this that the eigenvalues of (A) are all positive.

\item\label{exer:13.2.9}Solve the eigenvalue problem
$$
y''+2y'+y+\lambda y=0,\quad y(0)=0, \quad y(1)=0.
$$


\item\label{exer:13.2.10}Solve the eigenvalue problem
$$
y''+2y'+y+\lambda y=0,\quad y'(0)=0, \quad y'(1)=0.
$$

\exercisetext{
In Exercises~$\ref{exer:13.2.11}$--$\ref{exer:13.2.20}:$
\part{a} Determine whether $\lambda=0$ is an eigenvalue. If it is, find an
 associated eigenfunction.

\part{b}  Compute the negative eigenvalues with  errors not greater than
$5\times 10^{-8}$.  State the form of the associated eigenfunctions.


\part{c}  Compute the first four positive eigenvalues with  errors not
greater than
$5\times 10^{-8.}$  State the form of the associated eigenfunctions.}

 \item\label{exer:13.2.11}\Cex
$y''+\lambda y=0$, \quad  $y(0)+2y'(0)=0$,\quad $y(2)=0$

 \item\label{exer:13.2.12}\Cex
$y''+\lambda y=0$, \quad  $y'(0)=0$,\quad $y(1)-2y'(1)=0$


\item\label{exer:13.2.13}\Cex
$y''+\lambda y=0$, \quad  $y(0)-y'(0)=0$,\quad $y'(\pi)=0$


 \item\label{exer:13.2.14}\Cex
$y''+\lambda y=0$, \quad  $y(0)+2y'(0)=0$,\quad $y(\pi)=0$


\item\label{exer:13.2.15}\Cex
$y''+\lambda y=0$, \quad  $y'(0)=0$,\quad $y(2)-y'(2)=0$



 \item\label{exer:13.2.16}\Cex
$y''+\lambda y=0$, \quad  $y(0)+y'(0)=0$,\quad $y(2)+2y'(2)=0$


 \item\label{exer:13.2.17}\Cex
$y''+\lambda y=0$, \quad  $y(0)+2y'(0)=0$,\quad $y(3)-2y'(3)=0$



 \item\label{exer:13.2.18}\Cex
$y''+\lambda y=0$, \quad  $3y(0)+y'(0)=0$,\quad $3y(2)-2y'(2)=0$




 \item\label{exer:13.2.19}\Cex
$y''+\lambda y=0$, \quad  $y(0)+2y'(0)=0$,\quad $y(3)-y'(3)=0$


 \item\label{exer:13.2.20}\Cex
$y''+\lambda y=0$, \quad  $5y(0)+2y'(0)=0$,\quad $5y(1)-2y'(1)=0$




\item\label{exer:13.2.21}Find the first five eigenvalues of  the boundary value problem
$$
y''+2y'+y+\lambda y=0,\quad y(0)=0,\quad y'(1)=0
$$
with errors not greater than $5\times 10^{-8}$. State the form of the
associated eigenfunctions.


\exercisetext{In Exercises~\ref{exer:13.2.22}--\ref{exer:13.2.24} take it as given
that
$\{xe^{kx},xe^{-kx}\}$ and $\{x\cos kx, x\sin kx\}$ are fundamental sets of
solutions of
$$
x^{2}y''-2xy'+2y-k^{2}x^{2}y=0
$$
and
$$
x^{2}y''-2xy'+2y+k^{2}x^{2}y=0,
$$
respectively.}



\item\label{exer:13.2.22}Solve the eigenvalue problem  for

$$
x^{2}y''-2xy'+2y+\lambda x^{2}y=0, \quad y(1)=0,\quad y(2)=0.
$$



\item\label{exer:13.2.23}\Cex
Find the first five eigenvalues of
$$
x^{2}y''-2xy'+2y+\lambda x^{2}y=0, \quad y'(1)=0,\quad y(2)=0
$$
with errors no greater than $5\times 10^{-8}$. State the form of the
associated eienfunctions.


\item\label{exer:13.2.24}\Cex
Find the first five eigenvalues of
$$
x^{2}y''-2xy'+2y+\lambda x^{2}y=0, \quad y(1)=0,\quad y'(2)=0
$$
with errors no greater than $5\times 10^{-8}$. State the form of the
associated eienfunctions.





\item\label{exer:13.2.25}Consider the Sturm-Liouville problem
\begin{equation}   \tag{A}
y''+\lambda y=0,\quad y(0)=0,\quad y(L)+\delta y'(L)=0.
\end{equation}

\part{a} Show that (A)  can't have more than one negative eigenvalue,
and find the values of $\delta$ for which it has one.

\part{b} Find all values of $\delta$ such that $\lambda=0$ is an eigenvalue
of (A).

\part{c} Show that $\lambda=k^{2}$ with $k>0$ is an eigenvalue of (A)
if and only if
\begin{equation} \tag{B}
\tan kL=-\delta k.
\end{equation}

\part{d}
For
 $n=1$, $2$, \dots, let $y_{n}$  be an eigenfunction  associated with
$\lambda_{n}=k_{n}^{2}$. From Theorem~\ref{thmtype:13.2.4}, $y_{m}$  and
$y_{n}$ are orthogonal over $[0,L]$ if $m\ne n$. Verify this  directly.
\hint{Integrate by parts twice and use} (B).


\item\label{exer:13.2.26}Solve the Sturm-Liouville problem
$$
y''+\lambda y=0,\quad y(0)+\alpha y'(0)=0,\quad y(\pi)+\alpha y'(\pi)=0,
$$
where $\alpha\ne0$.



\item\label{exer:13.2.27}\Lex
Consider the Sturm-Liouville problem
\begin{equation}   \tag{A}
y''+\lambda y=0,\quad y(0)+\alpha y'(0)=0,\quad y(1)+(\alpha-1) y'(1)=0,
\end{equation}
where $0<\alpha<1$.

\part{a} Show that $\lambda=0$  is an eigenvalue of   (A),
and find an associated eigenfunction.


\part{b} Show that  (A)   has a negative eigenvalue,
and find the form of an associated eigenfunction.

\part{c}   Give a graphical argument to show that (A) has infinitely many
positive eigenvalues $\lambda_{1} < \lambda_{2}<\cdots
<\lambda_{n}<\cdots$, and state the form of the associated eigenfunctions.

\exercisetext{
Exercises~\ref{exer:13.2.28}--\ref{exer:13.2.30}  deal with the Sturm--Liouville
problem
\begin{equation} \tag{SL}
y''+\lambda y=0,\quad \alpha y(0)+\beta y'(0),\quad
\rho y(L)+\delta y'(L)=0,
\end{equation}
where $\alpha^{2}+\beta^{2}>0$  and  $\rho^{2}+\delta^{2}>0$.}


\item\label{exer:13.2.28}Show that $\lambda=0$ is an eigenvalue of (SL)  if and only if
$$
\alpha(\rho L+\delta)-\beta\rho =0.
$$


\item\label{exer:13.2.29}\Lex
The point of this exercise is that (SL) can't have more than two negative
eigenvalues.

\part{a} Show that $\lambda$ is a negative eigenvalue of (SL)
if and only if $\lambda=-k^{2}$, where $k$  is a positive solution of
$$
(\alpha\rho-\beta\delta k^{2})\sinh kL+k(\alpha\delta-\beta\rho)\cosh kL.
$$

\part{b}  Suppose $\alpha\delta-\beta\rho=0$.  Show that (SL)
has a negative eigenvalue if and only if $\alpha\rho$ and $\beta\delta$
are both nonzero. Find the negative eigenvalue and an associated
eigenfunction. \hint{Show that in this case $\rho=p\alpha$ and
$s=q\beta$, where $q\ne0$.}

\part{c}  Suppose $\beta\rho-\alpha\delta\ne0$. We know from
Section~11.1 that (SL) has no negative eigenvalues
if $\alpha\rho=0$ and $\beta\delta=0$. Assume that either $\alpha\rho\ne0$
or $\beta\delta\ne0$. Then we can rewrite (A) as
$$
\tanh kL=
\frac{k(\beta\rho-\alpha\delta)}{\alpha\rho-\beta\delta k^{2}}.
$$
By graphing both sides of this equation on the same axes (there are
several possibilities for the right side), show that it has at most
two positive solutions, so (SL) has at most two negative eigenvalues.


\item\label{exer:13.2.30}\Lex
The point of this exercise is that (SL) has infinitely many positive
eigenvalues $\lambda_{1}<\lambda_{2}<\cdots<\lambda_{n}<\cdots$, and
that $\lim_{n\to\infty}\lambda_{n}=\infty$.




\part{a} Show that $\lambda$ is a positive eigenvalue of (SL)
if and only if $\lambda=k^{2}$, where $k$ is a positive solution of
\begin{equation} \tag{A}
(\alpha\rho+\beta\delta k^{2})\sin kL+k(\alpha\delta-\beta\rho)\cos kL=0.
\end{equation}



\part{b} Suppose $\alpha\delta-\beta\rho=0$. Show that the positive
eigenvalues
 of (SL) are $\lambda_{n}=(n\pi/L)^{2}$, $n=1$, $2$, $3$, \dots.
\hint{Recall the hint in Exercise~\ref{exer:13.2.29}\part{b}.}

Now suppose $\alpha\delta-\beta\rho\ne0$. From
Section~11.1,  if $\alpha\rho=0$ and $\beta\delta=0$,
then (SL) has the eigenvalues
$$
\lambda_{n}=[(2n-1)\pi/2L]^{2},\quad  n=1,2,3, \dots
$$
(why?), so let's suppose  in addition that at least one of the products
$\alpha\rho$ and $\beta\delta$ is nonzero.  Then we can rewrite (A)
as
\begin{equation} \tag{B}
\tan kL=
\frac{k(\beta\rho-\alpha\delta)}
{\alpha\rho-\beta\delta k^{2}}.
\end{equation}
By graphing both sides of this equation on the same axes
(there are several possibilities for the right side), convince yourself of
the following:


\part{c} If $\beta\delta=0$, there's a positive integer $N$ such that (B)
has one solution $k_{n}$ in each of the intervals
\begin{equation} \tag{C}
\left((2n-1)\pi/L, (2n+1)\pi/L)\right),\quad n=N,N+1,N+2,\dots,
\end{equation}
and either
$$
\lim_{n\to\infty}\left(k_{n}-\frac{(2n-1)\pi}{2L}\right) =0
\text{\quad or\quad}
\lim_{n\to\infty}\left(k_{n}-\frac{(2n+1)\pi}{2L}\right)=0.
$$


\part{d}
If $\beta\delta\ne0$, there's a positive integer $N$  such that (B) has
one solution $k_{n}$ in each of the intervals (C) and
$$
\lim_{n\to\infty}\left(k_{n}-\frac{n\pi}{N}\right)=0.
$$

\item\label{exer:13.2.31}The following Sturm--Liouville problems are genera1izations of Problems
1--4 of Section~11.1.

\quad  Problem 1:   \quad $(p(x)y')'+\lambda r(x)y=0$,\quad
$y(a)=0$,\quad  $y(b)=0$


\quad  Problem 2:   \quad $(p(x)y')'+\lambda r(x)y=0$,\quad
$y'(a)=0$,\quad  $y'(b)=0$


\quad  Problem 3:   \quad $(p(x)y')'+\lambda r(x)y=0$,\quad
$y(a)=0$,\quad  $y'(b)=0$


\quad Problem 4:   \quad $(p(x)y')'+\lambda r(x)y=0$,\quad
$y'(a)=0$,\quad  $y(b)=0$

Prove: Problems~1--4  have no negative eigenvalues. Moreover, $\lambda=0$
is an eigenvalue of Problem~2 with associated eigenfunction $y_{0}=1$,
but $\lambda=0$ isn't  an eigenvalue of Problems 1, 3, and 4.
\hint{See the proof of Theorem~\ref{thmtype:11.1.1}.}


\item\label{exer:13.2.32}Show that the eigenvalues of the Sturm--Liouville
problem
$$
(p(x)y')'+\lambda r(x)y=0,\quad \alpha y(a)+ \beta y'(a)=0,\quad
\rho y(b)+\delta y'(b)
$$
are all positive if $\alpha\beta\le 0$, \,  $\rho\delta\ge 0$, and
$(\alpha\beta)^{2}+(\rho\delta)^{2}>0$.



\end{exerciselist}
